\scanpage{58}
Ainsi, il y a deux extensions, à ne pas
confondre, ayant comme noyau le
groupe de l'icosaèdre orienté
\[
\begin{aligned}
0\longrightarrow\operatorname{Aut}(\vec I)&\longrightarrow
 \operatorname{Aut}(\underbrace{\{\vec I,\vec I^{\unsure{\circ|o}}\}}_{\text{bicos. or.}})
 \mathrel{\unsure{-|\to}}\mathbb Z/2\mathbb Z\longrightarrow0\\
0\longrightarrow\operatorname{Aut}(\vec I)&\mathrel{-}
 \operatorname{Aut}(\{I,I^\circ\})\longrightarrow\mathbb Z/4\mathbb Z\longrightarrow0
\end{aligned}
\]

Il n'y a pas, je crois, d'hom \unsure{dans|de} un sens
ou de l'autre, induisant l'identité
sur $\operatorname{Aut}(\vec I)$. De haut en bas, le gén. de
$\mathbb Z/2\mathbb Z$ (qui définit l'autom ext. non trivial
de $\operatorname{Aut}(\vec I)$) devrait aller sur l'un des
deux gén. de $\mathbb Z/4\mathbb Z$, impossible !

Dans l'autre sens, un \deleted{\unsure{autom}.} $\mathbb Z/4\mathbb Z\longrightarrow\mathbb Z/2\mathbb Z$
doit être surjectif ($=$ \unsure{raison|raison.|va}) le noyau
\deleted{\unsure{doit être|devrait}} de $\operatorname{Aut}(I,I^{\prime})$ est un sous-groupe
d'ordre $2$ -- l'antipodisme ! Et en fait,
je vois maintenant que $\operatorname{Aut}(\vec I,\vec I^{\unsure{\prime|\circ}})$ est \inserted{iso}
\deleted{le quotient} au quotient de \deleted{\unsure{par}} $\operatorname{Aut}(I,I')$
par l'antipodisme. Mais cela semblerait
indiquer qu'il y aurait un foncteur
\[
\text{biicosaèdres (droits) non orienté}\longmapsto\substack{\text{biicosaèdre}\\\text{orienté}}
\]
et en effet, on passe par biicosaèdres
gauches, et biicosaèdres orientés via rev.
des orientations.
% archive sheet number: 57
